\label{app:bhd_for_ligo_sensing_and_control}
The sensors and actuators for controlling the various degrees of freedom that the HAM6 layout is responsible for is described in brief in the following sections.
\subsection{Longitudinal Control}
Longitudinal control refers to the sensing and actuation of lengths within the interferometer, i.e.\,position of optics along the optical axis. An example of longitudinal degree of freedom is the differential arm length. There are many longitudinal degrees of freedom within the interferometer and many of them are the lengths of cavities. The degrees of longitudinal degree of freedom that the HAM6 layout must provide sensing and control for a described in the next three sections.
\subsubsection{Homodyne Angle}
One method of controlling the homodyne angle is using the so called `RF 36` signal. LIGO uses 9\,MHz and 45\,MHz modulations on the light to detect various lengths within the interferometer that need to be controlled. The beat between these two will produce a 36,MHz signal. The homodyne angle determines the amplitude of this signal, and so by demodulating a photodiode measuring the light in reflection from the OMC at 36\,MHz the homodyne angle may be controlled. This is detailed in (Tengs Thesis ref). The actuation may be provided by OM0, BHDBS2 or by BHDM1. Alternatively, a dither signal at a few kHz on OM0 could provide can provide a signal corresponding to the homodyne angle.
\subsubsection{OMC Length}
Keeping the OMC length locked to the combined LO-AS beam can be done by actuating on the length of PZT mirrors in the OMC. A few kHz signal may be created by the modulation on a HRTS in either the LO or AS path, one of the HDS after BHDBS2, or event a dither on the other PZT within the OMC. These modulations will be measurable in transmission of the OMC. Alternatively, a PDH signal may be created by using the 9\,MHz signal may be measured in reflection from the OMC.
\subsubsection{Squeezing Angle to Measurement Phase}
The squeezing and measurement phase may be controlled by measuring the 3,MHz signal transmitted through the OMC. This signal originates in the AOM in the squeezer, and the squeezing angle can be controoled with the PZT in the the squeezer layout in HAM7.
\subsection{Alignment Control}
The alignment of beams within the interferometer is required to maximise the power of the laser the mode containing the gravitational wave signal. The alignment degrees of freedom that need to be sensed and controlled by the HAM6 layout are discussed in the next four sections.
\subsubsection{LO and AS alignment}
The alignment of the signal to the to the AS beam will be done with either OM0 or BHDM1, and BHDBS2. A 36\,MHz signal measured in reflection of the OMC can be used to determine the overlap of the two beams. By having a Gouy Phase telescope in reflection of the the OMC, a signal proportional to displacement and angle may be obtained. Alternatively, a dither at a few Hz can be applied to BHDM1 and BHDBS2, and this signal will appear in the light transmitted by the OMC.
\subsubsection{Alignment of the LO-AS beam to the OMC}
A dither applied to a HRTS in either the LO or AS path will produce a signal in transmission through the OMC. The mirrors OMxS and OMx2 will be used to align the LO beam into the OMC, and as these are separated by $\sim90\si{\degree}$ of Gouy phase, orthogonal control of the input beam alignment will be achievable. 
\subsubsection{Squeezer to AS beam}
The pick-off at OM0 will be used to align the squeezer beam to the AS beam. The squeezer light has 42\,MHz sidebands, and so a pair of wavefront sensors separated by $\sim90\si{\degree}$ of Gouy phase can be demodulated at this frequency to derive an alignment error signal. The squeezer beam will be aligned to the AS beam with optics between the filter cavity and the OFI. The signal and squeezer beam will be aligned onto these wavefront sensors using SR2 and a tip tilt in transmission of OM0.
\subsubsection{Differential Hard Mode}
The differential hard mode arises due to radiation pressure effects in the arm cavities. As the power in the cavity changes, the dynamics of the test mass changes. These dynamics result from the alignment of the cavity axis to the test masses due radiation pressure. This differential hard mode dynamic must be controlled. This signal will be detected with light that is transmitted through OM0, and actuated on with the test masses.